% Ficha de estudio 1 del Taller virtual ATA EPE: cinemática de cuerpos rígidos.
\documentclass[11pt,a4paper]{article}
% Preámbulo común de las fichas de estudio del Taller virtual ATA EPE.
\usepackage[spanish,es-noshorthands,es-tabla]{babel}
\usepackage[utf8]{inputenc}
\usepackage[T1]{fontenc}
\usepackage{lmodern}
\usepackage[margin=2.2cm]{geometry}
\usepackage[table]{xcolor}
\usepackage{amsmath,amssymb,bm}
\usepackage{tabularx,array,booktabs}
\usepackage{enumitem}
\usepackage{fontawesome5}
\usepackage[most]{tcolorbox}
\usepackage[hidelinks]{hyperref}
\usepackage{listings}

\definecolor{morado}{HTML}{2A1B7E}
\definecolor{amarillo}{HTML}{FFCF2E}
\definecolor{rojo}{HTML}{C4142C}
\definecolor{verde}{HTML}{1F9D55}
\definecolor{gris}{HTML}{5D5A73}

\newcolumntype{Y}{>{\raggedright\arraybackslash}X}
\renewcommand{\arraystretch}{1.25}
\setlist{itemsep=1pt,topsep=3pt}
\numberwithin{equation}{section}
\lstset{basicstyle=\ttfamily\small,columns=fullflexible,keepspaces=true,frame=single,rulecolor=\color{gris!40},
  literate={á}{{\'a}}1 {é}{{\'e}}1 {í}{{\'i}}1 {ó}{{\'o}}1 {ú}{{\'u}}1 {ñ}{{\~n}}1 {Ω}{{$\Omega$}}1
    {θ}{{$\theta$}}1 {λ}{{$\lambda$}}1 {Δ}{{$\Delta$}}1 {μ}{{$\mu$}}1 {·}{{$\cdot$}}1 {←}{{$\leftarrow$}}1 {§}{{\S}}1 {→}{{$\to$}}1 {≠}{{$\neq$}}1 {É}{{\'E}}1 {Á}{{\'A}}1 {∨}{{$\vee$}}1 {∧}{{$\wedge$}}1 {σ}{{$\sigma$}}1 {×}{{$\times$}}1 {≈}{{$\approx$}}1}

\newcommand{\rev}[1]{\widetilde{#1}}
\newcommand{\I}{\mathbf{I}}
\newcommand{\vect}[1]{\bm{#1}}
\newcommand{\A}{\mathcal{A}}

\newtcolorbox{meta}{colback=morado,colframe=morado,coltext=white,arc=3mm}
\newtcolorbox{lectura}{colback=morado!5,colframe=morado,fonttitle=\bfseries,title={\faBook~Lecturas}}
\newtcolorbox{ejercicio}[1]{colback=amarillo!10,colframe=amarillo!80!black,fonttitle=\bfseries,title={\faPencil*~Ejercicio #1}}
\newtcolorbox{control}{colback=rojo!5,colframe=rojo,fonttitle=\bfseries,title={\faQuestionCircle~Preguntas de control}}
\newtcolorbox{codigo}[1][De la ecuación al código]{colback=verde!6,colframe=verde,fonttitle=\bfseries,title={\faCode~#1}}
\newtcolorbox{cierre}{colback=verde!6,colframe=verde,fonttitle=\bfseries,title={\faFlagCheckered~Cierre del módulo}}
\newtcolorbox{nota}[1]{colback=amarillo!12,colframe=amarillo!80!black,fonttitle=\bfseries,title={\faLightbulb~#1}}

% \cabecera{número}{título}{subtítulo}
\newcommand{\cabecera}[3]{%
  \title{\textcolor{morado}{\textbf{Ficha #1. #2}}\\[3pt]\large Taller virtual ATA EPE · #3}
  \author{}\date{}}

\cabecera{1}{Movimientos rígidos, tornillos y velocidad}{Módulo 1: cinemática}

\begin{document}
\maketitle
\vspace{-2.5em}

\begin{meta}
\textbf{\textcolor{amarillo}{Meta del módulo.}} Describir cualquier movimiento de una pieza como un tornillo,
derivarlo en el tiempo para obtener velocidades, interpolar entre poses, llevar un agujero sobre otro y redondear una
pose a la retícula del kit. Cara pública: \textbf{armar} — tomar piezas, moverlas, girarlas e imantarlas.
\end{meta}

\begin{lectura}
\begin{enumerate}[leftmargin=*]
  \item \textbf{Lynch y Park, \emph{Modern Robotics}, cap.~3} (§3.1--3.3): movimientos rígidos, coordenadas
        exponenciales, \emph{twists} y tornillos. Es la versión matricial de lo que haremos en PGA: léelo en paralelo.
  \item \textbf{Dorst y De Keninck, \emph{May the Forque Be with You}}, secciones de motores y de cinemática
        (velocidad como bivector, marco del cuerpo y del espacio).
  \item Murray, Li y Sastry, \emph{A Mathematical Introduction to Robotic Manipulation}, cap.~2 (§2.1--2.4): el
        teorema de Chasles y la teoría de tornillos con rigor.
  \item Selig, \emph{Geometric Fundamentals of Robotics}, caps.~4 y 7: el grupo $SE(3)$, su álgebra de Lie y los
        tornillos; puente con el álgebra de Clifford.
  \item Opcional e histórico: Ball, \emph{A Treatise on the Theory of Screws}, caps.~1--2.
\end{enumerate}
\end{lectura}

% =================================================================================================
\section{El grupo de los movimientos rígidos}

Un movimiento rígido conserva distancias y orientación. Forman un grupo, $SE(3)$, de dimensión 6 (3 de giro y 3 de
traslación). Hay varias maneras de representarlo:

\begin{tabularx}{\linewidth}{@{}l c Y@{}}
\toprule
\textbf{Representación} & \textbf{Números} & \textbf{Comentario}\\
\midrule
Matriz homogénea $4\times4$ & 12 (16) & la de three.js; componer es multiplicar matrices\\
Cuaternión y vector & 7 & común en motores físicos; dos objetos separados\\
Cuaternión dual & 8 & un solo objeto; es el motor de PGA con otra notación\\
Motor de PGA & 8 & un solo objeto que además actúa igual sobre planos, rectas y puntos\\
\bottomrule
\end{tabularx}

\medskip
El motor tiene 8 componentes (grados 0, 2 y 4) pero solo 6 grados de libertad: cumple $M\rev M = 1$, que son dos
ecuaciones (la parte escalar y la de grado 4).

\begin{ejercicio}{1.1}
Escribe un motor general $M = a + b_{01}e_{01} + \dots + b_{23}e_{23} + c\,\I$ y desarrolla $M\rev M$. Identifica las
dos ecuaciones de normalización. ¿Qué significa geométricamente la de grado 4?
\end{ejercicio}

% =================================================================================================
\section{Chasles y el tornillo}

Todo movimiento rígido es un \textbf{tornillo}: un giro de ángulo $\phi$ alrededor de una recta $L$ más un avance $s$ a
lo largo de ella. Su \emph{paso} es $h = s/\phi$ (en el kit, el tornillo M4 real tiene $h = 0{,}7\,\mathrm{mm}/2\pi$).
En PGA (ficha 0, §4.2):
\begin{equation}
  M = \exp\!\Bigl(-\tfrac12(\phi + s\,\I)\,L\Bigr).
\end{equation}
Dos casos límite: $h = 0$ es un giro puro (una bisagra); $h = \infty$ es una traslación pura (un riel).

\begin{ejercicio}{1.2}
Una tuerca enroscada que gira una vuelta completa avanza 0{,}7\,mm. Escribe su motor como exponencial de un único
bivector. ¿Qué pasa con ese motor si la tuerca da diez vueltas? ¿Y media vuelta en sentido contrario?
\end{ejercicio}

% =================================================================================================
\section{Velocidad: derivar un motor}

Si una pieza tiene pose $M(t)$, derivar $M\rev M = 1$ muestra que $\dot M \rev M$ es un bivector. Se define la
velocidad en dos marcos:
\begin{equation}
  \dot M = -\tfrac12\,\Omega_e\,M = -\tfrac12\,M\,\Omega_c, \qquad \Omega_e = M\,\Omega_c\,\rev M .
\end{equation}
$\Omega_c$ (marco del cuerpo) es la velocidad «vista desde la pieza»; $\Omega_e$ (marco del espacio) es la vista desde
la mesa. La relación entre ambas es la \emph{adjunta}: el sándwich. Un punto $X$ de la pieza, en el mundo
$X' = MX\rev M$, se mueve con
\begin{equation}
  \dot X' = X'\times\Omega_e \quad\text{(salvo el signo de la convención)},
\end{equation}
donde $\times$ es el conmutador. La recta de $\Omega_e$ es el eje instantáneo de tornillo.

\begin{ejercicio}{1.3}
Deriva $X' = M X \rev M$ usando (3.1) y llega a (3.2). Fija el signo con la convención de la ficha 0: un rotor de
velocidad angular $+\omega$ alrededor de $+z$ debe mover el punto $(1,0,0)$ hacia $+y$.
\end{ejercicio}

\begin{ejercicio}{1.4}
Para $\Omega = \omega\,e_{12} + v\,e_{01}$ (giro en $z$ más velocidad en $x$, con la convención de la ficha 0),
encuentra el eje instantáneo y su paso. ¿Qué pasa cuando $\omega \to 0$?
\end{ejercicio}

% =================================================================================================
\section{Interpolar poses}

Para animar el imantado (la pieza «cae» suave a su lugar) se interpola entre dos poses $M_0$ y $M_1$ siguiendo el
tornillo que las une:
\begin{equation}
  M(t) = \exp\!\bigl(t\,\log(M_1\rev M_0)\bigr)\,M_0,\qquad t\in[0,1].
\end{equation}
Es la generalización a $SE(3)$ de la interpolación esférica de cuaterniones (\emph{slerp}): velocidad constante a lo
largo del tornillo, sin las deformaciones de interpolar matrices componente a componente.

% =================================================================================================
\section{Llevar un agujero sobre otro}

Para dos rectas normalizadas $A$ y $B$ (dos ejes de agujeros), el motor más corto que lleva $A$ sobre $B$ es la raíz
cuadrada del producto $BA^{-1} = -BA$:
\begin{equation}
  M_{A\to B} = \sqrt{-BA},\qquad \sqrt{N} = \frac{1 + N}{\lVert 1 + N\rVert}\quad\text{(si }N\neq-1\text{)}.
\end{equation}

\begin{ejercicio}{1.5}
Verifica que $\sqrt N$ así definida cumple $(\sqrt N)^2 = N$ para un rotor $N = \cos\alpha + L\sin\alpha$. ¿Qué falla
cuando $N = -1$? Muestra que $-BA = -1$ ocurre cuando $B$ es la misma recta que $A$ con la orientación opuesta (un
agujero «al revés») y propone cómo resolver ese caso.
\end{ejercicio}

\begin{ejercicio}{1.6}
Para alinear agujeros, el ángulo y la distancia entre dos rectas salen de $AB$ (plan, sección de PGA). Calcula $AB$ para
dos rectas paralelas al eje $z$ que pasan por $(0,0)$ y $(d,0)$. ¿En qué parte de $AB$ aparece $d$?
\end{ejercicio}

% =================================================================================================
\section{La retícula del kit}

Las piezas del kit se arman en una retícula de 10\,mm con giros de $90^\circ$. Las orientaciones posibles son las 24
rotaciones del cubo (el grupo octaédrico $O$). Imantar una pose $M$ es encontrar el elemento más cercano de
\begin{equation}
  \mathcal{G} = \{\,T_{\vect{n}}\,R_g \;:\; \vect{n}\in(1\,\mathrm{cm})\,\mathbb{Z}^3,\ g\in O\,\}.
\end{equation}
Se redondea por separado: primero el giro al elemento de $O$ más cercano (el de mayor $\lvert\langle R_g\rev R\rangle_0\rvert$),
luego la traslación a la retícula. Pero las piezas no se imantan a la mesa sino \emph{entre sí}: lo natural es buscar
el agujero más cercano de otra pieza y usar (5.1).

\begin{ejercicio}{1.7}
Construye las 24 rotaciones del cubo como rotores (pista: identidad; 9 giros de $90^\circ$, $180^\circ$ y
$270^\circ$ en $x$, $y$, $z$; 6 giros de $180^\circ$ en diagonales de caras; 8 de $120^\circ$ y $240^\circ$ en
diagonales del cubo). ¿Cuántos rotores distintos hay, si $R$ y $-R$ representan el mismo giro?
\end{ejercicio}

% =================================================================================================
\section{Elegir con el dedo: el rayo como recta}

El cursor en la pantalla define un rayo desde la cámara: en PGA es la recta $L = C\vee P$, con $C$ el centro de la
cámara y $P$ el punto del cursor en el plano cercano. Elegir una pieza es encontrar la primera forma que corta $L$;
arrastrarla sobre la mesa es seguir el punto $L\wedge p_{\text{mesa}}$, donde la recta corta el plano de la mesa.

\begin{codigo}
\begin{lstlisting}
rayo(cámara, x_pantalla, y_pantalla):
    C ← punto(centro de la cámara)
    P ← punto(desproyectar(x, y) en el plano cercano)
    devolver normalizar(vee(C, P))
puntoSobreMesa(L):
    devolver wedge(L, plano(0, 1, 0, 0))     # la mesa es y = 0
\end{lstlisting}
\end{codigo}

% =================================================================================================
\section{Puertos de las piezas}

Cada agujero es una recta en el sistema de la pieza. Hoy \texttt{herramientas/generar\_piezas.py} conoce los
agujeros pero no los exporta: hay que agregar al \texttt{catalogo.json} la lista de puertos de cada pieza (centro,
dirección, diámetro y largo) y convertirlos en rectas al cargar.

\begin{codigo}[Qué programar]
\begin{tabularx}{\linewidth}{@{}l Y@{}}
\texttt{nucleo/cinematica.js} & \texttt{velocidadDePunto}, \texttt{aEspacio}/\texttt{aCuerpo} (adjunta),
  \texttt{interpolar}, \texttt{raiz}, \texttt{motorEntreRectas}, \texttt{anguloYDistancia}(A, B),
  \texttt{rotacionesDelCubo}, \texttt{imantar}.\\
\texttt{herramientas/generar\_piezas.py} & Exportar los puertos al catálogo.\\
\texttt{kit/catalogo.js} & Cargar piezas con sus puertos como rectas.\\
\texttt{app/} & Escena, rayo del cursor, mover, girar de a $90^\circ$, imantar con animación.\\
\end{tabularx}
\end{codigo}

\begin{control}
\begin{enumerate}[leftmargin=*]
  \item ¿Por qué el motor tiene 8 números pero 6 grados de libertad?
  \item ¿Qué diferencia hay entre $\Omega_c$ y $\Omega_e$, y cuándo conviene cada una?
  \item ¿Qué es el eje instantáneo de tornillo y cómo se obtiene de $\Omega$?
  \item ¿Por qué interpolar con (4.1) es mejor que interpolar matrices componente a componente?
  \item ¿Por qué la raíz de un motor falla cuando $N = -1$, y cuándo ocurre al alinear dos agujeros?
  \item ¿Por qué las orientaciones del kit forman un grupo de 24 elementos?
\end{enumerate}
\end{control}

\begin{cierre}
Ejercicios 1.1--1.7 · preguntas de control · pruebas del módulo (rectas, interpolación, raíz, imantado) · cara
pública: armar la estructura de la balanza en menos de 2 minutos.
\end{cierre}

\end{document}
